\PassOptionsToPackage{unicode}{hyperref}
\documentclass[11pt,a4paper]{article}
\usepackage[T1]{fontenc}
\usepackage[utf8]{inputenc}
\usepackage{lmodern,microtype,xcolor}
\usepackage{tcolorbox}
\tcbuselibrary{skins,breakable,listings}
\usepackage{titlesec,fancyhdr,enumitem,booktabs,longtable,array,colortbl}
\usepackage[margin=2.4cm]{geometry}
\usepackage{listings}
\usepackage{hyperref}

\definecolor{javaOrange}{HTML}{E76F00}
\definecolor{javaDeep}{HTML}{BF5700}
\definecolor{javaBlue}{HTML}{1565C0}
\definecolor{javaTeal}{HTML}{00695C}
\definecolor{javaAmber}{HTML}{B45309}
\definecolor{javaInk}{HTML}{2B2140}
\definecolor{javaCodeBg}{HTML}{FFF8F0}
\definecolor{javaCodeInline}{HTML}{BF5700}
\definecolor{javaRule}{HTML}{FFD7B5}
\definecolor{javaQuoteBg}{HTML}{FFFAF5}
\color{javaInk}

\lstdefinestyle{java}{
  language=Java,
  basicstyle=\small\ttfamily\color{javaInk},
  keywordstyle=\color[HTML]{D73A49}\bfseries,
  stringstyle=\color[HTML]{032F62},
  commentstyle=\color[HTML]{6A737D}\itshape,
  breaklines=true,breakatwhitespace=true,
  tabsize=4,showstringspaces=false,
  morekeywords={var,record,String,Scanner,ArrayList,HashMap,System,Override,
    public,private,protected,class,interface,extends,implements,new,this,super,
    return,void,static,final,abstract,null,true,false,throws,throw,try,catch,finally},
}
\newtcblisting{javacode}{
  listing only,listing style=java,
  breakable,enhanced jigsaw,arc=4pt,boxrule=0pt,frame hidden,
  colback=javaCodeBg,borderline west={3pt}{0pt}{javaDeep},
  left=10pt,right=8pt,top=6pt,bottom=6pt,
  before skip=8pt,after skip=8pt,
}
\renewcommand{\texttt}[1]{{\color{javaCodeInline}\small\ttfamily #1}}
\renewenvironment{quote}
  {\begin{tcolorbox}[breakable,enhanced jigsaw,arc=4pt,boxrule=0pt,frame hidden,
    colback=javaQuoteBg,borderline west={3pt}{0pt}{javaOrange},
    left=12pt,right=10pt,top=8pt,bottom=8pt,before skip=8pt,after skip=8pt]\itshape}
  {\end{tcolorbox}}
\titleformat{\section}{\sffamily\Large\bfseries\color{javaOrange}}{\thesection}{0.7em}{}
  [{\vspace{2pt}\color{javaRule}\titlerule[1.2pt]}]
\titleformat{\subsection}{\sffamily\large\bfseries\color{javaDeep}}{\thesubsection}{0.6em}{}
\titleformat{\subsubsection}{\sffamily\normalsize\bfseries\color{javaTeal}}{\thesubsubsection}{0.5em}{}
\titleformat{\paragraph}[runin]{\sffamily\bfseries\color{javaAmber}}{}{0em}{}
\titlespacing*{\section}{0pt}{2.2ex plus 1ex minus .2ex}{1.2ex plus .2ex}
\setlength{\headheight}{24pt}
\pagestyle{fancy}\fancyhf{}
\renewcommand{\headrulewidth}{0.6pt}
\fancyhead[L]{\sffamily\small\color{javaDeep}Java Programming MOOC}
\fancyhead[R]{\sffamily\small\color{javaDeep}\leftmark}
\fancyfoot[C]{\sffamily\small\color{javaOrange}\thepage}
\renewcommand{\headrule}{\hbox to\headwidth{\color{javaRule}\leaders\hrule height \headrulewidth\hfill}}
\setlist[itemize]{leftmargin=1.4em,itemsep=2pt,topsep=3pt}
\setlist[enumerate]{leftmargin=1.6em,itemsep=2pt,topsep=3pt}
\renewcommand{\labelitemi}{\textcolor{javaOrange}{\textbullet}}
\renewcommand{\arraystretch}{1.25}
\arrayrulecolor{javaRule}
\hypersetup{colorlinks=true,linkcolor=javaDeep,urlcolor=javaBlue,citecolor=javaTeal}
\makeatletter
\newcommand{\subtitle}[1]{\gdef\@subtitle{#1}}
\providecommand{\@subtitle}{}
\def\@maketitle{%
  \newpage\null\vskip 2em%
  \begin{center}%
    {\sffamily\bfseries\Huge\color{javaOrange}\@title\par}\vskip 1em%
    {\sffamily\Large\color{javaDeep}\@subtitle\par}\vskip 1.2em%
    {\color{javaRule}\rule{0.6\linewidth}{1.4pt}\par}\vskip 1em%
    {\sffamily\large\color{javaInk}\@author\par}\vskip 0.4em%
    {\sffamily\small\color{javaAmber}\@date\par}%
  \end{center}\par\vskip 2em}
\makeatother

\title{Programação em Java}
\subtitle{Lição 5: Orientação a Objetos --- Aprofundamento}
\author{Java Programming MOOC · Universidade de Helsinque --- tradução para o português}
\date{2026}

\begin{document}
\maketitle
{\hypersetup{linkcolor=javaDeep}\setcounter{tocdepth}{3}\tableofcontents}
\newpage

\section{Lição 5: Orientação a Objetos --- Aprofundamento}

\subsection{Objetivos de aprendizagem}

Ao final desta lição, você será capaz de:

\begin{itemize}
\item Criar múltiplos construtores e métodos sobrecarregados
\item Entender a diferença entre variáveis primitivas e de referência
\item Comparar objetos com o método \texttt{equals()}
\item Usar objetos como variáveis de instância (composição)
\item Compreender o que é \texttt{null} e como evitar \texttt{NullPointerException}
\end{itemize}

\subsection{Sobrecarga de construtores}

Uma classe pode ter \textbf{vários construtores}, desde que difiram no número ou tipo de parâmetros. Isso se chama \textbf{sobrecarga} (\textit{overloading}).

\begin{javacode}
public class Pessoa {
    private String nome;
    private int idade;

    // Construtor com apenas o nome
    public Pessoa(String nome) {
        this.nome = nome;
        this.idade = 0;
    }

    // Construtor com nome e idade
    public Pessoa(String nome, int idade) {
        this.nome = nome;
        this.idade = idade;
    }
}
\end{javacode}

Uso:

\begin{javacode}
Pessoa alice = new Pessoa("Alice");       // idade = 0
Pessoa bruno = new Pessoa("Bruno", 30);  // idade = 30
\end{javacode}

\subsubsection{Evitando repetição com \texttt{this()}}

Um construtor pode chamar outro usando \texttt{this()}. Essa chamada deve ser a \textbf{primeira instrução}:

\begin{javacode}
public Pessoa(String nome) {
    this(nome, 0); // chama o construtor de dois parâmetros
}

public Pessoa(String nome, int idade) {
    this.nome = nome;
    this.idade = idade;
}
\end{javacode}

\subsection{Sobrecarga de métodos}

Da mesma forma, podemos ter \textbf{múltiplos métodos com o mesmo nome}, diferindo apenas nos parâmetros:

\begin{javacode}
public void crescer() {
    this.idade++;
}

public void crescer(int anos) {
    this.idade += anos;
}
\end{javacode}

Uso:

\begin{javacode}
alice.crescer();     // envelhece 1 ano
alice.crescer(5);    // envelhece 5 anos
\end{javacode}

Versão sem repetição (usando \texttt{this.metodo()}):

\begin{javacode}
public void crescer() {
    this.crescer(1);
}

public void crescer(int anos) {
    this.idade += anos;
}
\end{javacode}

\subsection{Variáveis primitivas vs.\ variáveis de referência}

\subsubsection{Primitivas}

Os tipos \texttt{int}, \texttt{double}, \texttt{boolean}, \texttt{char} etc.\ armazenam diretamente o valor na memória:

\begin{javacode}
int a = 5;
int b = a; // b recebe uma COPIA do valor 5

b = 10;
System.out.println(a); // 5 - a nao mudou
System.out.println(b); // 10
\end{javacode}

\subsubsection{Referências}

Objetos (como \texttt{String}, \texttt{ArrayList}, ou qualquer classe criada) armazenam uma \textbf{referência} ao objeto na memória:

\begin{javacode}
Pessoa p1 = new Pessoa("Ana");
Pessoa p2 = p1; // p2 aponta para o MESMO objeto que p1

p2.crescer();
System.out.println(p1.getIdade()); // 1 - p1 tambem mudou!
\end{javacode}

Para ter \textbf{dois objetos independentes}, crie cada um separadamente:

\begin{javacode}
Pessoa p1 = new Pessoa("Ana");
Pessoa p2 = new Pessoa("Ana"); // objeto diferente
\end{javacode}

\subsubsection{O valor \texttt{null}}

Uma variável de referência que não aponta para nenhum objeto tem o valor \texttt{null}:

\begin{javacode}
Pessoa p = null;
System.out.println(p); // null

p.getNome(); // ERRO: NullPointerException!
\end{javacode}

Sempre verifique se uma referência é \texttt{null} antes de usá-la:

\begin{javacode}
if (p != null) {
    System.out.println(p.getNome());
}
\end{javacode}

\subsection{Comparando objetos com \texttt{equals()}}

O operador \texttt{==} compara \textbf{referências} (se apontam para o mesmo objeto), não os \textbf{conteúdos}:

\begin{javacode}
Pessoa p1 = new Pessoa("Ana", 25);
Pessoa p2 = new Pessoa("Ana", 25);

System.out.println(p1 == p2); // false - sao objetos diferentes na memoria
\end{javacode}

Para comparar pelo \textbf{conteúdo}, implemente o método \texttt{equals()}:

\begin{javacode}
@Override
public boolean equals(Object obj) {
    if (this == obj) return true;           // mesmo objeto
    if (!(obj instanceof Pessoa)) return false; // tipo diferente

    Pessoa outra = (Pessoa) obj;
    return this.nome.equals(outra.nome) && this.idade == outra.idade;
}
\end{javacode}

\begin{javacode}
System.out.println(p1.equals(p2)); // true - mesmos atributos
\end{javacode}

\subsection{Objetos como variáveis de instância (composição)}

Um objeto pode \textbf{conter} outros objetos como variáveis de instância. Isso se chama \textbf{composição}:

\begin{javacode}
public class Endereco {
    private String rua;
    private String cidade;
    private String cep;

    public Endereco(String rua, String cidade, String cep) {
        this.rua = rua;
        this.cidade = cidade;
        this.cep = cep;
    }

    @Override
    public String toString() {
        return this.rua + ", " + this.cidade + " - " + this.cep;
    }
}
\end{javacode}

\begin{javacode}
public class Pessoa {
    private String nome;
    private Endereco endereco; // composicao!

    public Pessoa(String nome, Endereco endereco) {
        this.nome = nome;
        this.endereco = endereco;
    }

    @Override
    public String toString() {
        return this.nome + "\n  " + this.endereco;
    }
}
\end{javacode}

Uso:

\begin{javacode}
Endereco end = new Endereco("Rua das Flores, 100", "Sao Paulo", "01310-100");
Pessoa p = new Pessoa("Maria", end);
System.out.println(p);
\end{javacode}

\textbf{Saída:}
\begin{javacode}
Maria
  Rua das Flores, 100, Sao Paulo - 01310-100
\end{javacode}

\subsection{Exemplo completo: Datas simples}

\begin{javacode}
public class Data {
    private int dia;
    private int mes;
    private int ano;

    public Data(int dia, int mes, int ano) {
        this.dia = dia;
        this.mes = mes;
        this.ano = ano;
    }

    public boolean ehAnterior(Data outra) {
        if (this.ano != outra.ano) return this.ano < outra.ano;
        if (this.mes != outra.mes) return this.mes < outra.mes;
        return this.dia < outra.dia;
    }

    @Override
    public boolean equals(Object obj) {
        if (!(obj instanceof Data)) return false;
        Data outra = (Data) obj;
        return this.dia == outra.dia && this.mes == outra.mes && this.ano == outra.ano;
    }

    @Override
    public String toString() {
        return String.format("%02d/%02d/%04d", this.dia, this.mes, this.ano);
    }
}
\end{javacode}

\begin{javacode}
Data natal = new Data(25, 12, 2026);
Data anoNovo = new Data(1, 1, 2027);

System.out.println(natal.ehAnterior(anoNovo)); // true
System.out.println(natal.equals(anoNovo));      // false
System.out.println(natal);                      // 25/12/2026
\end{javacode}

\subsection{Continuando o aprendizado de OOP}

À medida que os programas crescem, a OOP torna-se ainda mais valiosa. Na próxima parte, veremos como separar a interface do usuário da lógica do programa e como escrever testes.

\subsubsection{Dicas práticas}

\begin{itemize}
\item \textbf{Responsabilidade única}: cada classe deve ter um único propósito bem definido.
\item \textbf{Evite objetos gigantes}: se uma classe ficou muito grande, considere dividi-la.
\item \textbf{Prefira composição}: em vez de tentar resolver tudo numa classe só, componha classes menores.
\end{itemize}

\subsection{Resumo}

\begin{itemize}
\item \textbf{Sobrecarga}: múltiplos construtores/métodos com o mesmo nome mas parâmetros diferentes.
\item \texttt{this()} chama outro construtor da mesma classe (deve ser a primeira linha).
\item \textbf{Primitivas} copiam valores; \textbf{referências} copiam ponteiros para objetos.
\item \texttt{null} representa ausência de objeto; acessar métodos em \texttt{null} causa \texttt{NullPointerException}.
\item \texttt{==} compara referências; \texttt{equals()} compara conteúdo (implemente-o para suas classes).
\item \textbf{Composição}: um objeto pode conter outros objetos como atributos.
\end{itemize}

\subsection{Exercícios}

\paragraph{Exercício 1 --- Sobrecarga de construtores}

Crie a classe \texttt{Livro} com os atributos \texttt{titulo}, \texttt{autor} e \texttt{anoPublicacao}. Implemente:
\begin{itemize}
\item Construtor com os três atributos
\item Construtor apenas com \texttt{titulo} e \texttt{autor} (ano = 0)
\item Construtor apenas com \texttt{titulo} (autor = "Desconhecido", ano = 0)
\item \texttt{toString()} e \texttt{equals()} (igualdade por título e autor)
\end{itemize}

\paragraph{Exercício 2 --- Referências vs.\ primitivas}

Escreva um programa que demonstre a diferença entre copiar um \texttt{int} e copiar uma referência a um objeto. Mostre que alterar a cópia de um \texttt{int} não afeta o original, mas alterar via uma referência copiada afeta o objeto original.

\paragraph{Exercício 3 --- Ponto no plano}

Crie a classe \texttt{Ponto} com \texttt{x} e \texttt{y} (double). Implemente:
\begin{itemize}
\item \texttt{distanciaAte(Ponto outro)} usando \texttt{Math.sqrt}
\item \texttt{equals(Object obj)} com precisão de 0.001
\item \texttt{toString()} no formato \texttt{(x, y)}
\end{itemize}

Crie a classe \texttt{Circulo} com um \texttt{Ponto} centro e um raio \texttt{double}. Implemente:
\begin{itemize}
\item \texttt{contemPonto(Ponto p)} --- verifica se o ponto está dentro do círculo
\item \texttt{toString()}
\end{itemize}

\paragraph{Exercício 4 --- Null safety}

Crie um método \texttt{imprimirPessoa(Pessoa p)} que:
\begin{itemize}
\item Se \texttt{p} for \texttt{null}, exibe \texttt{"Pessoa não encontrada."}
\item Caso contrário, exibe as informações da pessoa
\end{itemize}
Teste chamando o método com \texttt{null} e com um objeto válido.

\end{document}
